\emph{Telescope: 150/750 Newton. For this task, we recommend using the 10 mm
eyepiece and the 2$\times$ Barlow lens.}

\textbf{Observation of Saturn} \hfill available time: 10 minutes

\begin{itemize}
\item In the upper box, the circle represents the disk of Saturn and the
  horizontal line is the E--W direction on the sky. Pay attention to the
  direction of North (see top right corner).

  Mark the position of Titan by a cross.
\item The smaller box on the bottom right corner of the first box is for
  drawing the rings of Saturn. Again the circle represents the disk of
  Saturn.

  Draw the rings of Saturn in this box with the correct size and
  orientation. Both the outer and inner edges of the ring are necessary, no
  faint ring details or gaps are needed. Keep the orientation of the image
  the same as the orientation in the upper box.
\item Estimate the angular distance (in arcsec) and position angle (in
  degrees) of Titan relative to the center of Saturn. You may do your
  calculations besides the answer.
\end{itemize}

Apparent major axis of the ring: $43''$

% FIGURE NEEDED: answer boxes -- large circle (disk of Saturn) with E-W line
% and North direction marked, plus smaller inset box for drawing the rings

\medskip
\emph{Alternative task (used in case of bad sky conditions at low
altitudes): \textbf{Epsilon Lyrae}}

\begin{itemize}
\item Find $\epsilon$ Lyr, and make a drawing of the field of view (with the
  object and other stars) with the 10 mm eyepiece. Label the directions
  North and East by two arrows and mark them as `N' and `E'.
\item Estimate the angular distance between the wide pair
  ($\epsilon_1$--$\epsilon_2$), and estimate the position angle of the same
  pair.
\item Increase the magnification with the 2$\times$ Barlow lens to be able
  to resolve and separate the two close pairs. Estimate the angle (in
  degrees to the nearest integer) subtended by the two close pairs relative
  to each other (the enclosed angle of the two lines going through the two
  narrow pairs). Do not give any PA, only the relative angle of the two
  close pairs. No drawing is needed.
\end{itemize}

If you cannot find $\epsilon$ Lyr, the assistant can point to it for you,
but only after 5 minutes. In this case you will lose the marks for pointing
the telescope to the object.
